\section{Momenta：量产、客户价值与“鲶鱼”}

Waymo 这一章解释了系统如何被拆解，Momenta 这一章则解释系统如何被客户检验。对高继扬来说，从 Waymo 到 Momenta，不是从好公司跳到另一家好公司，而是从工程师训练场进入量产交付现场。

如果说 Waymo 是工程师天堂，Momenta 则是量产战场。高继扬选择 Momenta，是因为他想从自动驾驶系统走向产品、客户和公司经营，尤其想理解一个 demo 级系统如何变成车企可交付的产品。

\begin{figure}[H]
\centering
\includegraphics[width=0.96\textwidth]{waymo-momenta-contrast.png}
\caption{Waymo 与 Momenta：两种极端组织训练。}
\end{figure}

\begin{knowledgebox}{读图：两个组织分别训练什么}
左侧 Waymo 训练系统工程、infra、代码质量和长期战略执行；右侧 Momenta 训练交付、客户沟通、组织调整和商业价值循环。一个让人学会系统怎样 work，另一个让人学会系统怎样被客户使用、付费和倒逼迭代。
\end{knowledgebox}

\subsection{为什么量产是数据路线}

Momenta 的核心逻辑是：自动驾驶最终需要大量真实数据。如果只靠自营车队，覆盖城市和场景太慢；如果把辅助驾驶或泊车等能力装到量产车上，先为客户创造价值，再通过量产车获得数据，就能形成商业驱动的数据循环。这里的关键不是“卖给车企”本身，而是数据获取不再是纯成本，而是被客户价值支付。

\begin{importantbox}{数据与商业价值循环}
量产不是单纯交付项目，而是让产品进入真实车、真实客户和真实道路，从而把数据采集变成商业驱动行为。客户价值越明确，数据回流越可持续。
\end{importantbox}

\subsection{鲶鱼角色}

前面讲 Momenta 的组织转型，这里把镜头拉回高继扬本人。所谓鲶鱼，不是一个性格标签，而是一种在高压交付环境里被反复锻炼出来的系统介入方式。

访谈标题里的“鲶鱼”，指高继扬在 Momenta 被灵活地放到不同模块：感知、定位、泊车、infra、规控、NOA 量产。他自己的总结是：快速进入不熟悉领域，用固定方法论理解系统，拆解任务，做人事匹配，监控反馈，反馈好就扩大，反馈不好就收缩调整。

\begin{knowledgebox}{鲶鱼方法论}
鲶鱼不是到处搅动，而是把系统变得更会反应。它的动作包括：进入陌生领域、快速建立地图、拆任务、配人、设指标、看反馈、调组织。
\end{knowledgebox}

\subsection{本章小结}

Momenta 给高继扬的是客户价值和组织压力的训练。它让他看到，一个技术团队若要变成产品公司，就必须经历交付、淘汰、调整和量产的洗礼。

